\chapter{Problem statement}
\label{ch:problem_statement}

\textbf{Motivation.} The application areas for robotics continues to expand, even in household contexts. Users may not always be able to receive training on how to operate robotic systems, as formal training can be very time consuming for individuals and for companies. Intuitive forms of teleoperation aim to solve this, by reading an operator's movements/gestures and mapping these to corresponding commands (such as setpoints) for the robot systems. This introduces a new challenge; mapping human limb kinematics to corresponding kinematics for robotic limbs, which often aren't anthropomorphic. There are many mathematical methods (models and transformations) that solve this problem, but few generalised implementations \cite{bib:teleoperation_continuous_space}\cite{bib:teleoperation_continuous_space_2}. Configuring and implementing these mathematical models remains a challenge and the implementation is often case-specific \cite{bib:teleoperation_intuitive_using_MediaPipeHands}\cite{bib:teleoperation_surgery}. This project aims to develop a software framework that simplifies the configuration and implementation of one or many of these methods, thus reducing the time required to develop and re-configure teleoperation systems. In addition to developing the framework, this project will test it in a real-world use-case, to demonstrate its effectiveness. Since cybersecurity becomes a major concern in the context of teleoperated systems, this project briefly raises and discusses concerns about remote cyber-physical systems, using the real-world implementation as an example.

\textbf{Context.} This project aims to develop a software framework that simplifies the configuration and implementation of one or many of the mapping methods. This will allow more time and resources to be invested into cutting-edge technologies of teleoperation project, and make these kinda of projects more accessible to hobbyists who may not be proficient in kinematics. Optimistically, if the framework becomes widely adopted, it may introduce new standardisations to teleoperation, inviting experts to optimise the framework, and develop surrounding/cooperative frameworks/solutions.

\textbf{Technical challenge.} Some of the tehcnical challenges of this project are: (i) to develop the mapping framework to be easy/intuitive, (ii) versatile, (iii) and efficient. Challenges relating to the real-world testing include: (iv) implementing a computer vision system to detect the user's hand positions, (v) interfacing with the framework, (vi) wirelessly and securely communicating setpoints to a Braccio arm, (vii) and making the Braccio arm move to the given setpoints.

\textbf{Limitations.} Depending on the computer vision implementation, the user's workspace may be limited, such as needing a static plain background. The scope of this project does not include optimising the real-world demonstration.
